% sections2/sec6_contraction.tex
\section{张量缩并}
\label{sec:contraction}

设 $V$ 为域 $K$ 上的有限维向量空间.求值配对
\[
  V \times V^\vee \to K, \quad (v, \varphi) \mapsto \varphi(v)
\]
是 $K$-双线性的，因此诱导出唯一的 $K$-线性映射
\begin{equation}
  c \colon V \otimes_K V^\vee \to K, \quad c(v \otimes \varphi) = \varphi(v).
\end{equation}
该映射称为\newterm{tensor-contraction}{张量缩并}{tensor contraction}.

在基底与对偶基底下的表现如下：设 $\{e_1, \dots, e_n\}$ 为 $V$ 的一组基，$\{e^1, \dots, e^n\}$ 为 $V^\vee$ 的对偶基.若 $v = \sum v^i e_i$ 且 $w = \sum w_j e^j$，则
\begin{equation}
  c(v \otimes w) = \sum_{i, j} v^i w_j e^j(e_i) = \sum_{i} v^i w_i. \label{eq:6.1}
\end{equation}
在同构 $V \otimes_K V^\vee \cong \End_K(V)$ 下，缩并运算 $c$ 恰好对应于矩阵的\textbf{迹}（Trace）：
\[
  \Tr(T) = \sum_{i} T^i_i.
\]

更一般地，对于混合张量空间 $V^{\otimes k} \otimes (V^\vee)^{\otimes \ell}$（$k, \ell \geq 1$），选取第 $r$ 个逆变分量与第 $s$ 个协变分量进行对偶求值，诱导出一般张量缩并映射：
\begin{equation}
  c_{r, s} \colon V^{\otimes k} \otimes (V^\vee)^{\otimes \ell} \to V^{\otimes (k-1)} \otimes (V^\vee)^{\otimes (\ell-1)},
\end{equation}
在 Einstein 求和约定下，缩并简记为将一对上下重复的相同指标进行求和，例如：
\[
  T^{i_1 \dots i_k}_{j_1 \dots j_\ell} \mapsto T^{i_1 \dots m \dots i_k}_{j_1 \dots m \dots j_\ell}.
\]

当空间 $V$ 上赋予了非退化对称双线性型（度规张量）$g \colon V \xrightarrow{\cong} V^\vee$ 时，我们可以利用 $g$ 实现\termcn{降指标}（$v_j = \sum_i g_{ij} v^i$），或利用其逆 $g^{-1}$ 实现\termcn{升指标}（$v^j = \sum_i g^{ij} v_i$），从而将同类张量（如 $V \otimes_K V$）通过复合映射
\begin{equation}
  V \otimes_K V \xrightarrow{\id \otimes g} V \otimes_K V^\vee \xrightarrow{c} K
\end{equation}
进行度规缩并：$T^{ij} \mapsto \sum_{i, j} g_{ij} T^{ij}$.这是黎曼几何与广义相对论中所有标量曲率及场方程运算的数学基石.
